\documentclass[tikz,border=4pt]{standalone}
\usepackage{ctex}
\usepackage{pgfplots}
\pgfplotsset{compat=1.18}
\usetikzlibrary{calc}

\begin{document}
% part12/03 频率收敛到概率的折线图（抛硬币正面频率）
\begin{tikzpicture}
\begin{axis}[
  axis lines=left,
  xlabel={试验次数 $n$},
  ylabel={正面出现频率},
  xmin=0, xmax=1050,
  ymin=0, ymax=1.05,
  ytick={0,0.25,0.5,0.75,1.0},
  yticklabels={$0$, $0.25$, $\mathbf{0.5}$, $0.75$, $1.0$},
  xtick={0,200,400,600,800,1000},
  width=9cm, height=6cm,
  title={频率收敛到概率（抛硬币实验）},
  title style={font=\small},
  xlabel style={font=\small},
  ylabel style={font=\small},
  tick label style={font=\small},
  extra y ticks={0.5},
  extra y tick style={
    grid=major,
    grid style={red, dashed, thick},
    tick style={draw=none},
    ticklabel style={color=red, font=\small},
  },
]
  % 模拟频率收敛折线（数据近似模拟大数定律，非精确随机）
  \addplot[blue, thick] coordinates {
    (10,  0.70)
    (30,  0.60)
    (50,  0.56)
    (100, 0.53)
    (150, 0.54)
    (200, 0.52)
    (300, 0.51)
    (400, 0.505)
    (500, 0.50)
    (600, 0.502)
    (700, 0.499)
    (800, 0.500)
    (900, 0.501)
    (1000, 0.500)
  };
  % 理论概率线 p=0.5
  \addplot[red, thick, dashed, domain=0:1000] {0.5};
  \node[red, right, font=\small] at (axis cs:1010,0.5) {$P=0.5$};
  \node[blue, font=\footnotesize, right] at (axis cs:200,0.68) {频率};
\end{axis}
\end{tikzpicture}
\end{document}
